\section{Notation}
\label{sec:background}


Our work assumes a multi-agent task that is commonly described as a stochastic game $G$, specified by a tuple $G = {\langle}S, U, P, r, Z, O, n, \gamma{\rangle}$. Here $n$ agents, $a \in A \equiv \{1,...,n\}$, choose actions, $u^a \in U$, and $s \in S$ is the state of the environment. The joint action $\mathbf{u} \in \mathbf{U} \equiv U^{n}$ leads to a state transition based on the transition function $P(s'|s,\mathbf{u}): S \times \mathbf{U} \times S \rightarrow [0,1]$. The reward functions $r^a(s,\mathbf{u}): S \times \mathbf{U} \rightarrow \mathbb{R}$ specify the rewards and $\gamma \in [0,1)$ is the discount factor.

We further define the discounted future return from time $t$ onward as $R^a_t = \sum_{l=0}^\infty \gamma^l r^a_{t+l}$ for each agent, $a$. In a naive approach, each agent maximises its total discounted return in expectation separately. This can be done with policy gradient methods \citep{sutton1999policy} such as REINFORCE~\citep{williams1992simple}. Policy gradient methods update an agent's policy, parameterised by $\vct{\theta}^a$, by performing gradient ascent on an estimate of the expected discounted total reward $\exs{}{R^a_0}$. 

By convention, bold lowercase letters denote column vectors.



